% This is text associated with the open textbook "Open Electromagnetics"
% See immediately below for author, date
% This work is licensed under the Creative Commons Attribution-ShareAlike 4.0 International License. (CC BY-SA 4.0) 
% To view a copy of the license, visit http://creativecommons.org/licenses/by-sa/4.0/

%ORB TITLE Power Flow on Transmission Lines
%ORB AUTHOR 0        # S.W. Ellingson, Virginia Tech (ellingson@vt.edu)
%ORB DATE 20180802   # yyyymmdd
%ORB ABSTRACT Power Flow on Transmission Lines

\label{m0090_Power_Flow_on_Transmission_Lines} % <-- Must match name of this file and module directory name.  Don't mess with this!

It is often important to know the power associated with a wave on a transmission line.
The power of the waves incident upon, reflected by, and absorbed by a load are each of interest.
In this section we shall work out expressions for these powers and consider some implications in terms of the voltage reflection coefficient ($\Gamma$) and standing wave ratio (SWR).
  
Let's begin by considering a lossless transmission line that is oriented along the $z$ axis.
The time-average power associated with a sinusoidal wave having potential $v(z,t)$ and current $i(z,t)$ is
%
\begin{equation}
P_{av}(z) \triangleq \frac{1}{T}\int_{t_0}^{t_0+T}{ v(z,t)~i(z,t)~dt }
\end{equation} 
%
where $T\triangleq 2\pi/f$ is one period of the wave and $t_0$ is the start time for the integration.
Since the time-average power of a sinusoidal signal does not change with time, $t_0$ may be set equal to zero without loss of generality. 

Let us now calculate the power of a wave incident from $z<0$ on a load impedance $Z_L$ at $z=0$.  
We may express the associated potential and current as follows:
%
\begin{equation}
v^+(z,t) = \left|V_0^+\right|\cos\left(\omega t -\beta z +\phi\right)
\end{equation}
%
\begin{equation}
i^+(z,t) = \frac{\left|V_0^+\right|}{Z_0}\cos\left(\omega t -\beta z +\phi\right)
\end{equation}
%
And so the associated time-average power is
%
\begin{equation}
\begin{split}
P_{av}^+(z) &= \frac{1}{T}\int_0^T{ v^+(z,t)~i^+(z,t)~dt } \\
            &= \frac{\left|V_0^+\right|^2}{Z_0} \cdot \frac{1}{T}\int_0^T{ \cos^2\left(\omega t -\beta z +\phi\right) ~dt } 
\end{split}
\label{m0090_ePav1}
\end{equation}
%
Employing a well-known trigonometric identity:
\begin{equation}
\cos^2\theta = \frac{1}{2} + \frac{1}{2}\cos 2\theta 
\end{equation}
we may rewrite the integrand as follows
%
\begin{equation}
\cos^2\left(\omega t -\beta z +\phi\right) = \frac{1}{2} + \frac{1}{2}\cos\left(2\left[\omega t -\beta z +\phi\right]\right) 
\end{equation}
%
Then integrating over both sides of this quantity
%
\begin{equation}
\int_0^T{ \cos^2\left(\omega t -\beta z +\phi\right) ~dt } = \frac{T}{2} + 0 
\end{equation}
%
The second term of the integral is zero because it is the integral of cosine over two complete periods.
Subsequently, we see that the position dependence (here, the dependence on $z$) is eliminated.  
In other words, the power associated with the incident wave is the same for all $z<0$, as expected.
Substituting into Equation~\ref{m0090_ePav1} we obtain:
%
\begin{equation}
\boxed{
P_{av}^+ = \frac{\left|V_0^+\right|^2}{2Z_0} 
}
\label{m0090_ePav+}
\end{equation}
%
This is the time-average power associated with the incident wave, measured at any point $z<0$ along the line.

%ORB HIGHLIGHT
\begin{mdframed}[backgroundcolor=yellow!20,nobreak=true] 
Equation~\ref{m0090_ePav+} gives the time-average power 
\index{transmission line!time-average power}
associated with a wave traveling in a single direction along a lossless transmission line.
\end{mdframed}

Using precisely the same procedure, we find that the power associated with the \emph{reflected} wave is
%
\begin{equation}
P_{av}^- = \frac{\left|\Gamma V_0^+\right|^2}{2Z_0} = \left|\Gamma\right|^2 \frac{\left|V_0^+\right|^2}{2Z_0}
\end{equation}
%
or simply
%
\begin{equation}
\boxed{
P_{av}^- = \left|\Gamma\right|^2 P_{av}^+
}
\label{m0090_ePav-}
\end{equation}

%ORB HIGHLIGHT
\begin{mdframed}[backgroundcolor=yellow!20,nobreak=true] 
Equation~\ref{m0090_ePav-} gives the time-average power associated with the wave reflected from an impedance mismatch.
\end{mdframed}

Now, what is the power $P_L$ delivered to the load impedance $Z_L$?
The simplest way to calculate this power is to use the principle of \emph{conservation of power}.
Applied to the present problem, this principle asserts that the power incident on the load must equal the power reflected plus the power absorbed; i.e.,
%
\begin{equation}
P_{av}^+ = P_{av}^- + P_L
\end{equation}
%
Applying the previous equations we obtain:
%
\begin{equation}
\boxed{
P_L = \left(1-\left|\Gamma\right|^2\right) P_{av}^+
}
\label{m0090_ePL}
\end{equation}

%ORB HIGHLIGHT
\begin{mdframed}[backgroundcolor=yellow!20,nobreak=true] 
Equations~\ref{m0090_ePL} gives the time-average power transferred to a load impedance, and is equal to the difference between the powers of the incident and reflected waves. 
\end{mdframed}

%ORB EXAMPLE
~\\ \begin{example} 
\begin{mdframed}[backgroundcolor=brown!20,linecolor=brown!20]\raggedright
How important is it to match $50~\Omega$ to $75~\Omega$?
\\~\\
Two impedances which commonly appear in radio engineering are $50~\Omega$ and $75~\Omega$.  
It is not uncommon to find that it is necessary to connect a transmission line having a $50~\Omega$ characteristic impedance to a device, circuit, or system having a $75~\Omega$ input impedance, or vice-versa.  If no attempt is made to match these impedances, what fraction of the power will be delivered to the termination, and what fraction of power will be reflected?  What is the SWR?

{\bf Solution.} The voltage reflection coefficient going from $50~\Omega$ transmission line to a $75~\Omega$ load is
%
\begin{equation}
\Gamma = \frac{75-50}{75+50} = +0.2
\nonumber
\end{equation}
%
The fraction of power reflected is $\left|\Gamma\right|^2 = 0.04$, which is 4\%.
The fraction of power transmitted is $1-\left|\Gamma\right|^2$, which is 96\%. 
Going from a $50~\Omega$ transmission line to a $75~\Omega$ termination changes only the sign of $\Gamma$, and therefore, the fractions of reflected and transmitted power remain 4\% and 96\%, respectively.
In either case (from Section~\ref{m0081_Standing_Wave_Ratio}):
%
\begin{equation}
\mbox{SWR} = \frac{1+\left|\Gamma\right|}{1-\left|\Gamma\right|} = 1.5
\nonumber
\end{equation}
%
This is often acceptable, but may not be good enough in some particular applications.
Suffice it to say that it is not necessarily required to use an impedance matching device to connect $50~\Omega$ to $75~\Omega$ devices.
\end{mdframed}\end{example}

